\section{Posterior estimates at a stable field}\label{sec:consequences}

The residual rules now yield the two posterior estimates used in
Section~\ref{sec:observation}. We first select a finite collection of
small-residual regions, then compare the primary iterates with the exact
root. Throughout, $J$ satisfies the required fixed-length word diagnostics,
and $h_0$ satisfies the whole stable-field implication through residual
threshold $\rho_0\sqrt n$. Denote the unique zero of $F_{h_0}$ by $r$ and
set $t_*=\tanh r$, $q_*=\av{t_*^2}$. All expectations and forms here are
at $\mu_{h_0}$, and $R_l=m^{l-1}-m^l$ denotes the residual from
Section~\ref{sec:residuals}.

\subsection{Selecting a small-residual region}
The primary recursion need not converge. The finite-depth fact we need is
that two consecutive residuals are small on most of the Gibbs mass, with
a separate bound after square reweighting. Smooth cutoffs make that selection
compatible with discrete differentiation.
\begin{lemma}\label{lem:residual-selection}
For each sufficiently small fixed $\rho>0$, there is a fixed depth $d$ and
functions $C_k$, $2\le k\le d$, with the following properties:
\begin{enumerate}
\item $0\le C_k\le1$, and $C_k$ is supported where
$\max(\norm{R_k},\norm{R_{k-1}})\le\rho\sqrt n$. It equals one where both
norms are at most $\rho\sqrt n/2$, and $\norm{dC_k}\le C_\rho/\sqrt n$
pointwise.
\item Set $\widehat C_k=C_k\prod_{i=2}^{k-1}(1-C_i)$ and
$D=1-\sum_{k=2}^d\widehat C_k$. Then $0\le D\le1$ and
\[
 \Prob_{\mu_{h_0}}(D>0)\le\frac{C_\rho}{n}.
\]
For $N=\E f^2>0$ and $\nu=f^2\mu_{h_0}/N$,
\[
 \Prob_\nu(D>0)\le\frac{C_\rho}{\sqrt n}
                   \left(1+\frac{\D_{h_0}(f)}{N}\right).
\]
\end{enumerate}
All constants are uniform over the field and matrix under the stated
diagnostics. The depth and the required diagnostic length may depend on
$\rho$.
\end{lemma}

\begin{proof}
We use the scalar residual pairing $S_l=R_l^{\mathsf T}m^l/\sqrt n$ to
control the total squared increment. First, the primary derivative bounds
give $\norm{dS_l}\le C_l$ pointwise. In the product rule, the linear
terms are controlled by the operator norms of $\nabla R_l$ and
$\nabla m^l$, multiplied by $\norm{m^l}/\sqrt n$ and
$\norm{R_l}/\sqrt n$. The cross term for coordinate $i$ has size at most
$2\norm{d_i R_l}\norm{d_i m^l}/\sqrt n$; summing its square uses the
bounded column norms and gives the same scale.

Apply Lemma~\ref{lem:residual} with target $m^l/\sqrt n$ and test $G=S_l$.
It gives
\[
 \E S_l^2\le C_l\bigl(\E S_l^2+C_l\bigr)^{1/2},
 \qquad\text{hence}\qquad \E S_l^2\le C_l'.
\]
For the square-density version use $H_0=S_l/\sqrt n$, which is bounded with
bounded difference norm. Dividing the resulting estimate by $N\sqrt n$
gives
\begin{equation}\label{eq:weighted-S}
 \E_\nu\left(\frac{S_l}{\sqrt n}\right)^2
 \le\frac{C_l}{\sqrt n}\left(1+\frac{\D_{h_0}(f)}{N}\right).
\end{equation}
We now use the squared-increment identity
\eqref{eq:residual-telescoping}. Choose a fixed $d$ with
$\lfloor d/2\rfloor\rho^2/4>2$, followed by $\eps>0$ with $2d\eps<1$.
Whenever $\max_{l\le d}|S_l|/\sqrt n\le\eps$, telescoping gives
$\sum_{l=1}^d\norm{R_l}^2/n<2$. If every consecutive pair contained a
residual larger than $\rho\sqrt n/2$, the disjoint pairs
$(1,2),(3,4),\ldots$ would force the sum to exceed this bound. At least
one pair therefore meets the required threshold.

To implement the selection smoothly, apply a scalar cutoff equal to one
on $[0,1/4]$ and zero on $[1,\infty)$ to each of
$\norm{R_k}^2/(\rho^2n)$ and $\norm{R_{k-1}}^2/(\rho^2n)$, and multiply
the results to define $C_k$. The primary operator and column bounds,
with the product rule, give
$\norm{d(\norm{R_l}^2/n)}\le C_l/\sqrt n$, hence the claimed derivative
bound. Finally $D=\prod_{k=2}^d(1-C_k)$: positive deficit implies
$\max_{l\le d}|S_l|/\sqrt n>\eps$. The preceding two moment bounds and
a finite union bound give the ordinary and square-reweighted estimates.
\end{proof}

\subsection{A directional covariance bound}
On each selected region, the implicit test vector converts the residual
identity into a comparison with $t_*$. Its auxiliary scalar error is
estimated first; only then do we test in an arbitrary Euclidean direction.
\begin{proposition}\label{prop:directional}
There is a constant $C$ such that, at every field satisfying the above
stable-field implication,
\begin{equation}\label{eq:directional}
 \norm{\Cov_{\mu_{h_0}}(G,x)}
 \le C\bigl(\Var_{\mu_{h_0}}(G)+\D_{h_0}(G)\bigr)^{1/2}
\end{equation}
for every $G$. The constant is uniform in $n,h_0$, and $G$ on the specified
matrix events; only a fixed finite diagnostic length is required.
\end{proposition}

\begin{proof}
Fix an index $k\ge2$ and work on the small-residual region
\eqref{eq:small-region}, with the buffered implicit construction
\eqref{eq:implicit}. Abbreviate $m=m^k$, $p=m-t_*$, and put
\[
 a=j(1-b_{k-1}-q_*),\qquad
 L_* =\sqrt n\bigl(\chi-(1-q_*)\bigr).
\]
Here $\chi=\av{A_i}$ and $A_i=\Phi(h_i^k;r_i,a)$ are defined globally by
\eqref{eq:phi-data}, even though the implicit vector is only needed locally.
Thus $|L_*|\le C\sqrt n$ and $\norm{dL_*}\le C_k$ pointwise, by the primary
derivative bounds and $\norm{db_{k-1}}=O(n^{-1/2})$.

For either of the two allowed sources $w_0$, the following pointwise identity
holds on the buffered region:
\begin{align}\label{eq:implicit-comparison}
 &\bigl(h^k-h_0-(J-jb_{k-1}I)m\bigr)^{\mathsf T}w
       -j(b_k-b_{k-1})nc_1 \notag\\
 &\quad=\sum_i\bigl[(h_i^k-r_i-am_i)-a\partial_k\bigr]w_{0,i}
       +j(1-q_*-\chi)p^{\mathsf T}w-jc_1R_k^{\mathsf T}p.
\end{align}
Here the zeroth-order term in the first sum acts by multiplication on
$w_{0,i}$. For the verification, subtract the root equation to write the
vector on the left as $h^k-r-am-Jp+j(1-q_*)p$, then substitute
$Jw=y^{\rm imp}+j\chi w+jc_1m^{k-1}$. Apply \eqref{eq:phi-stein} to
the contribution $A(y^{\rm imp}-jc_1t_*)$, holding its auxiliary arguments
and parameters fixed, and use $nc_1=\sum_i\partial_k w_i$. The remaining
scalar terms cancel by
\[
 p^{\mathsf T}(m^{k-1}+t_*)+n(b_k-b_{k-1})
 =n(1-b_{k-1}-q_*)+R_k^{\mathsf T}p,
\]
which follows by expanding $p=m-t_*$ and $R_k=m^{k-1}-m$.

Lemma~\ref{lem:terminal-residual} controls the tested left-hand side of
\eqref{eq:implicit-comparison}. The last term on the right also has tested
expectation at most $C\norm{G}_*$: apply Lemma~\ref{lem:residual} to target
$p/\sqrt n$ and multiplier $\sqrt n c_1$. The latter is bounded with bounded
difference norm on the buffer. Meanwhile the size estimates for the implicit
solution and Lemma~\ref{lem:root} give
\begin{equation}\label{eq:middle-error}
 \abs{j(1-q_*-\chi)p^{\mathsf T}w}\le C_0\rho|L_*|.
\end{equation}
The source size and stability margins determine $C_0$ independently of
$k$ and the selection depth. This lets us choose $\rho$ first so that
$C_0\rho<1/2$ and all stability and flip-buffer conditions hold. The
selection depth in Lemma~\ref{lem:residual-selection} is chosen afterward.
This order will permit absorption of the scalar error without a
self-dependent choice of depth.

First control the averaged coefficient by taking
$w_{0,i}=\partial_{r_0}\Phi(h_i^k;r_i,a)/\sqrt n$.
The second identity in \eqref{eq:implicit-source-actions} makes the first
sum in \eqref{eq:implicit-comparison} equal to $L_*$.
Test with $G=C_k^2L_*$. The cutoff and difference bounds give
\[
 \norm{C_k^2L_*}_*\le C+\norm{C_kL_*}_{L^2(\mu_{h_0})}.
\]
Indeed its difference vector is bounded pointwise: $|L_*|=O(\sqrt n)$
multiplies $\norm{d(C_k^2)}=O(n^{-1/2})$, and $\norm{dL_*}=O(1)$.
Absorbing \eqref{eq:middle-error} now yields
\begin{equation}\label{eq:susceptibility-average}
 \norm{C_kL_*}_{L^2(\mu_{h_0})}\le C_k'.
\end{equation}

Now test a direction by taking $w_0=Ae$, where $e$ is any deterministic
vector with $\norm e\le1$. By the first identity in
\eqref{eq:implicit-source-actions}, the first sum in
\eqref{eq:implicit-comparison} is $p^{\mathsf T}e$.
Test with $\widehat C_kG$. The residual estimates control the left side
and the last term by $C\norm G_*$; the middle term is bounded by
\eqref{eq:susceptibility-average} and Cauchy--Schwarz, since
$0\le\widehat C_k\le C_k$. Thus
\[
 \abs{\E[\widehat C_kG(m^k-t_*)^{\mathsf T}e]}\le C\norm G_*.
\]
Summing Lemma~\ref{lem:residual} with the constant target $e$ over
$l=1,\ldots,k$ gives the same bound for
$\E[\widehat C_kG(x-m^k)^{\mathsf T}e]$.
Sum over $k$. The deficit has contribution at most
\[
 2\sqrt n\,\norm G_{L^2(\mu_{h_0})}
       \Prob_{\mu_{h_0}}(D>0)^{1/2}\le C\norm G_*.
\]
Thus $|\E[G(x-t_*)^{\mathsf T}e]|\le C\norm G_*$ uniformly over
$\norm e\le1$. Centering $G$ and taking the supremum over this seed
class gives \eqref{eq:directional}. The uncentered estimate also records
the ordinary-mean consequence: take $G=1$ to obtain
$\norm{\E_{\mu_{h_0}}x-t_*}\le C$. This norm is the unnormalized
Euclidean norm. It compares the mean itself with $t_*$; no
approximation of the derivative of the mean map is used here.
The local construction required $n$ above a fixed threshold. In the
remaining dimensions, the elementary bounds
$\norm{\Cov(G,x)}^2\le n\Var(G)$ and
$\norm{\E_{\mu_{h_0}}x-t_*}\le2\sqrt n$ extend both conclusions after
enlarging $C$.
\end{proof}

\subsection{Means under square reweighting}
For the change of observation law we need a different estimate, valid for
arbitrary square densities. It uses only ordinary recipes and trades a
small fixed leading error for a constant depending on that tolerance.
\begin{proposition}\label{prop:weighted-mean}
For every sufficiently small fixed $\delta>0$, there is a constant $C_\delta$
such that, for $N=\E_{\mu_{h_0}}f^2>0$ and $\nu=f^2\mu_{h_0}/N$,
\begin{equation}\label{eq:weighted-mean}
 \frac{\norm{\E_\nu x-t_*}}{\sqrt n}
 \le C\delta+\frac{C_\delta}{\sqrt n}
               \left(1+\frac{\D_{h_0}(f)}{N}\right).
\end{equation}
The leading constant $C$ is independent of $\delta$ and of the selection
depth. This estimate requires only ordinary word diagnostics, of a fixed
length that may depend on $\delta$.
\end{proposition}

\begin{proof}
Select the cutoffs with a fixed threshold $\rho\le\delta$ small enough
for Lemma~\ref{lem:root}. On each cutoff support, the residual-to-root
comparison gives
\[
 \norm{m^k-t_*}\le C\delta\sqrt n.
\]
Here the constant comes from the recurrence and the inverse bound in
Lemma~\ref{lem:root}, and is independent of $k$.
For $\norm e\le1$, sum the square-density version of
Lemma~\ref{lem:residual} with target $e$ through $k$, using multiplier
$\widehat C_k$. After division by $N$, this bounds
\[
 \abs{\E_\nu[\widehat C_k(x-m^k)^{\mathsf T}e]}
 \le C_\delta\left(1+\frac{\D_{h_0}(f)}{N}\right).
\]
The cutoff weights sum to at most one, so all root-comparison terms
together cost at most $C\delta\sqrt n$. The leftover mass contributes
at most $2\sqrt n\Prob_\nu(D>0)$, which is controlled by the second
part of Lemma~\ref{lem:residual-selection}. Add the terms, divide by
$\sqrt n$, and take the supremum over $e$ to obtain
\eqref{eq:weighted-mean}. The leading $C$ comes from root stability;
only $C_\delta$ and the finite diagnostic length depend on the selected
tolerance. Any dimensions below a fixed threshold required by this
construction are covered by the elementary bound
$\norm{\E_\nu x-t_*}/\sqrt n\le2$ after increasing $C_\delta$; this
leaves the leading constant unchanged. The argument has used no implicit
test vector.
\end{proof}
